\chapter{Diagramme utilisateur et parcours}
\label{chap:utilisateur}

\section{Cas d'utilisation}

\begin{figure}[htbp]
\centering
\begin{tikzpicture}[x=1mm,y=1mm]
  \node[actor] (user) at (0,-36) {Photographe};
  \node[component,text width=52mm] (open) at (75,0) {Ouvrir une image et son dossier};
  \node[component,text width=52mm] (browse) at (75,-18) {Parcourir et choisir une vignette};
  \node[component,text width=52mm] (adjust) at (75,-36) {Choisir le DCP et régler l'image};
  \node[component,text width=52mm] (meta) at (75,-54) {Noter, libeller et marquer};
  \node[component,text width=52mm] (batch) at (75,-72) {Sélectionner et appliquer en lot};
  \node[component,text width=52mm] (export) at (75,-90) {Exporter l'image active};
  \draw[thick,iprgray] (user.east) -- (30,-36);
  \draw[thick,iprgray] (30,0) -- (30,-90);
  \foreach \target in {open,browse,adjust,meta,batch,export}
    \draw[flow] (30,0 |- \target.west) -- (\target.west);
\end{tikzpicture}
\caption{Actions existantes. L'application de réglages en lot n'est pas un export par lot.}
\end{figure}

\section{Parcours nominal}

\begin{enumerate}
  \item L'utilisateur choisit \emph{Ouvrir}. Il peut sélectionner directement une image ; \IPR{} ouvre alors le dossier qui la contient et active l'image choisie.
  \item Le ruban inférieur affiche les vignettes. Une image sélectionnée devient l'image de travail ; le panneau d'information expose ses propriétés.
  \item Sur un RAW, le module Développement RAW / Profil expose les DCP compatibles.
  Le statut distingue normalisation demandée et effectivement appliquée.
  \item Les commandes de réglage modifient l'état de \texttt{ImageEditor}. Le champ numérique et le curseur représentent la même valeur et doivent rester synchronisés.
  \item Après l'arrêt de l'interaction, le rendu est mis à jour puis les réglages sont écrits dans le sidecar. Un changement d'image restaure ses propres réglages ou les valeurs neutres en l'absence de sidecar.
  \item L'utilisateur peut cocher plusieurs vignettes et appliquer les réglages de l'image active à la sélection. Chaque image reçoit alors son sidecar propre.
  \item L'export traite l'image active en pleine résolution. La file d'export
  d'une sélection entière reste à implémenter.
\end{enumerate}

\section{Exigences d'ergonomie}

Le panneau des réglages est fixé au bord droit et ne doit pas masquer l'image. Les barres de défilement et curseurs sont volontairement discrets, mais les champs numériques restent éditables afin de permettre une saisie précise. Les contrôles réinitialisent exclusivement leur réglage cible.
